\documentclass[../main.tex]{subfiles}

\begin{document}

\chapter{Introduction}\label{ch:introduction}

In today's increasingly digital world, where most online services require users to create an account to 
access them, users are forced to create and memorise a large number of different credentials, or are even 
forcibly encouraged to reuse old passwords on multiple services.
These online profiles are incredibly sensitive because they are usually linked to the personal identity 
of the user, such as their email address, phone number or even official government documentation like a 
personal ID.
More often than not, this profile is protected by a single password created by the user.
This dominant method of password-based authentication, which is responsible for verifying and authenticating 
users online, poses a significant risk to companies and other large organisations.
One way in which malicious attackers take advantage of this vulnerable authentication method is by using 
large databases of leaked user credentials (often including email addresses and corresponding passwords) 
to authenticate at Service Providers (SPs) as the affected user.
This attack method is called credential stuffing.
Verizon reports that credential abuse was the leading initial attack vector in 22\% of the 12,195 confirmed 
data breaches in 2025~\cite{VerizonDataBreach}.

To reduce the number of places where credentials are created and stored, the \acrfull{fim} was developed.
FIM allows users to authenticate on multiple SPs using the same credentials, which 
are provided by an \acrfull{idp}. 
FIM consolidates authentication behind a single trusted authority called the IdP.
When a user attempts to gain access on the \acrfull{sp} in FIM, the user's browser is redirected 
to the IdP.
The user authenticates with the IdP and the IdP communicates with the SP to guarantee that the user is 
who they say they are.

Shibboleth, developed by Internet2, is one of the most widely deployed open-source implementations of 
this model in the research and education (R\&E) sector~\cite{Shibboleth}. 
Under this model, a user authenticates with their organisation's IdP once.
The IdP then issues cryptographically signed SAML assertions that SPs accept 
as proof of identity.
This architecture is widely deployed in higher education and research federation such 
as DFN-AAI~\cite{DFNAAI}, 
where a single IdP serves dozens of SPs simultaneously.

This architectural property introduces a critical consequence: the IdP becomes a single point of trust 
for the entire federation~\cite{IdPSinglePointOfTrust}. 
A compromised, misconfigured or vulnerable IdP does not affect one application.
It undermines the security of every relying SP.
A notable example of this vulnerability is the Sunburst supply chain attack. 
After gaining initial access, the attackers obtained the trusted SAML signing certificate and used it to 
forge authentication assertions, granting unauthorised access to all connected services without triggering 
a single failed login~\cite{FireEyeSSC,MicrosoftBlogSSC}.
Beyond targeted attacks, the SAML protocol itself is susceptible to well-documented vulnerabilities, 
such as \acrfull{xsw}~\cite{VulnerSAMLCloud, BreakingSAML}, Response Injection, Assertion Artifact Leakage 
and IdP Spoofing~\cite{SecAnalysisSSO}.
These vulnerabilities enable attackers to forge or manipulate security claims accepted by the SP as 
authoritative once the signature has been validated.

Despite this elevated risk profile, systematic security monitoring at the application layer of Shibboleth 
IdP deployments remains largely absent. 
Existing approaches such as AAIEye~\cite{MonitoringAAI} focus primarily on operational availability, 
leaving security-relevant aspects unaddressed. 
The research literature reflects a comparable gap: a survey of security analyses in federated identity 
management finds Shibboleth to be among the least examined protocols, and formal analyses of SAML remain 
scarce despite its deployment across government, enterprise and education~\cite{FIMSurvey}.

This thesis addresses this gap by proposing a security-oriented monitoring system that operates directly 
within the Shibboleth application layer. 
Analogous to a network Looking Glass, which provides transparency into the internal routing state 
of a network node, the proposed system acts as a Security Looking Glass for the Shibboleth IdP. 
It makes application-layer security parameters and authentication workflows visible that are otherwise 
inaccessible from outside the application. 
Rather than testing the deployment actively, the system observes it, and thereby provides federation 
operators with a continuous view of the security posture of their authentication infrastructure.

Three research questions guide this work:
\begin{itemize}
    \item RQ1 Which security-relevant data can be collected at the application layer of a Shibboleth IdP, and through which mechanisms are these data accessible?
    \item RQ2 How can these data be correlated in order to detect configuration weaknesses and anomalies in the authentication process?
    \item RQ3 What measurable overhead does the resulting system impose on the monitored IdP?
\end{itemize}

The remainder of this thesis is structured as follows. 
\begin{itemize}
    \item Chapter \ref{ch:background} introduces the technical foundations of FIM, SAML and the Shibboleth IdP. 
    \item Chapter 3 positions this work within existing research on application-layer intrusion detection and delimits it from network- and host-based monitoring systems. 
    \item Chapter 4 presents the architecture of the proposed system, and 
    \item Chapter 5 describes its implementation devided in layers. 
    \item Chapter 6 evaluates the system with respect to the three research questions above. 
    \item Chapter 7 summarises the findings and outlines directions for future work.
\end{itemize}

\end{document}
